% Proposed-method section for the UC-RCF manuscript.
% Requires: amsmath, amssymb, booktabs, adjustbox.
%
% IMPORTANT---AUTHOR-ONLY PRE-SUBMISSION ALIGNMENT CHECKS (not rendered):
% 1. The current notebook sets BASE_VARIANT from the globally ranked ablation.
%    For the primary endpoint, either fix the residual-scale variant a priori or
%    select it inside every outer asset fold, then regenerate all headline results.
% 2. causal_roll currently applies bfill() to the first undefined samples. Replace
%    this with a past-only initialization before claiming strictly causal features.
% 3. The calibration-score "leave-self-out" correction in fit_variants has a
%    one-count discrepancy. Use rank/n after removing self, then rerun thresholds.
% 4. Replace the tautological donor-exclusion assertions in Cell 10 with explicit
%    set-intersection checks before calling the protocol leakage-resistant.
% 5. fit_variants currently counts missing channel p-values as non-firing in the
%    fraction score. Divide by the number of observed monitored channels, as
%    defined below, and regenerate the affected score and alarm results.

\section{Proposed method}
\label{sec:proposed_method}

\subsection{Overview and deployment setting}

The proposed uncertainty-calibrated residual conformal framework (UC-RCF)
addresses a deployment setting in which a target turbine has an available
healthy historical partition, but alarm calibration and operating-policy
selection must transfer across assets. This distinction is important: UC-RCF
is asset-adaptive rather than zero-shot. A normal-behavior model is estimated
from the designated healthy history of each target turbine, whereas the
case-level false-alarm behavior is calibrated using verified-normal prediction
periods from other turbines. The resulting separation prevents the prediction
period of a held-out asset from influencing its operating point.

Let $c\in\mathcal{C}$ index a labeled monitoring case, $a(c)$ its turbine, and
$f(c)$ its wind farm. Each case contains a chronologically ordered training
partition $\mathcal{T}^{\mathrm{tr}}_c$ and prediction partition
$\mathcal{T}^{\mathrm{pr}}_c$. At time $t$, the SCADA record is divided into an
operating vector $\mathbf{x}_{ct}\in\mathbb{R}^{d_c}$ and a vector of condition
responses $\mathbf{y}_{ct}\in\mathbb{R}^{m_c}$. Only rows satisfying the
dataset-defined normal operating states, $s_{ct}\in\{0,2\}$, are admitted to
normal-behavior fitting. Event annotations and case labels are not used to fit
the regression model or construct channel-wise conformal scores; they are used
only for development-fold policy selection and final evaluation.

UC-RCF comprises six stages: (i) farm-specific channel taxonomy and nonlinear
operating-feature construction; (ii) embargoed blocked cross-fitting of a
multi-output normal-behavior model; (iii) target screening and uncertainty
standardization of cross-fitted residuals; (iv) channel-wise conformal evidence
aggregation; (v) persistence-sensitive criticality accumulation; and (vi)
cross-asset calibration nested within leave-one-asset-out evaluation.

\subsection{Channel taxonomy and nonlinear operating representation}

The channel taxonomy is constructed from the CARE feature metadata and is
fixed independently of the event labels. Only average-value SCADA channels are
used. Wind speed, active and reactive power, rotor and generator speed,
ambient or nacelle temperature, pitch, and yaw-related variables form the
operating set. Counters are excluded, angular variables are encoded separately,
and all remaining nonconstant measurements form candidate condition-response
channels. This separation prevents a response channel from simultaneously
entering the model as an explanatory driver.

For each case, missing operating values are imputed using medians calculated
only from normal training rows. Every continuous driver is subsequently
standardized using the corresponding training mean and standard deviation and
clipped to $[-8,8]$. Denoting the standardized driver vector by
$\mathbf{z}_{ct}$, the nonlinear design is
\begin{equation}
\boldsymbol{\phi}_{ct}=
\left[
1,
\mathbf{z}_{ct}^{\top},
(\mathbf{z}_{ct}\odot\mathbf{z}_{ct})^{\top},
\left\{\mathcal{R}_{w}(\mathbf{z}_{ct,\mathcal{J}})^{\top}\right\}_{w\in\{6,36,144\}},
\left\{z_{ctj}z_{ctk}\right\}_{j<k;\,j,k\in\mathcal{J}},
\sin(\boldsymbol{\theta}_{ct})^{\top},
\cos(\boldsymbol{\theta}_{ct})^{\top}
\right]^{\top},
\label{eq:nonlinear_design}
\end{equation}
where $\mathcal{J}$ contains at most six core operating drivers and
$\mathcal{R}_{w}(\cdot)$ is a past-only rolling mean. At the nominal 10-min
sampling interval, $w\in\{6,36,144\}$ represents 1~h, 6~h, and 24~h operating
contexts, respectively. Equation~\eqref{eq:nonlinear_design} captures smooth
nonlinearity, interactions among dominant load variables, multiscale operating
history, and the periodic geometry of angular measurements without requiring a
high-capacity black-box model.

Candidate response channels are retained only when they have finite training
mean, non-negligible variance, and more than 50\% observed normal-training
values. Each retained response is standardized using its normal-training mean
and standard deviation. Missing standardized responses in the shared
multi-output regression solve are assigned their training-mean value (zero in
standardized units); unobserved prediction responses are excluded from the
subsequent evidence aggregation.

\subsection{Embargoed cross-fitted normal-behavior model}

Let $\mathbf{\Phi}_c$ be the design matrix formed by stacking
$\boldsymbol{\phi}_{ct}^{\top}$ over the normal training rows and let
$\mathbf{Y}_c$ contain the standardized candidate responses. UC-RCF estimates a
shared multi-output ridge model,
\begin{equation}
\widehat{\mathbf{B}}_{c,\lambda}
=\underset{\mathbf{B}}{\operatorname*{arg\,min}}
\left\|\mathbf{Y}_c-\mathbf{\Phi}_c\mathbf{B}\right\|_F^2
+\lambda\left\|\mathbf{B}\right\|_F^2,
\label{eq:ridge_objective}
\end{equation}
which provides one coherent operating representation for all response channels.
If $\mathbf{\Phi}_c=\mathbf{U}\operatorname{diag}(\mathbf{s})\mathbf{V}^{\top}$,
the solution is evaluated efficiently as
\begin{equation}
\widehat{\mathbf{B}}_{c,\lambda}
=\mathbf{V}\operatorname{diag}\!\left(
\frac{s_{\ell}}{s_{\ell}^{2}+\lambda}
\right)\mathbf{U}^{\top}\mathbf{Y}_c,
\label{eq:ridge_svd}
\end{equation}
allowing the full regularization path to share a single economy-size singular
value decomposition.

Random cross-validation is inappropriate for serially correlated SCADA data.
The normal training sequence is therefore divided into contiguous 432-sample
blocks (3~days at 10-min resolution), assigned round-robin to $K=5$ folds. When
a fold is held out, an embargo of 144 samples (24~h) is removed on both sides of
each held block before model fitting. The candidate grid is
$\Lambda=\{10^{-2},10^{-1},10^{0},\ldots,10^{5}\}$. For response $j$, let
$R^{2}_{cj}(\lambda)$ denote the coefficient of determination computed from
its embargoed out-of-fold predictions. A single case-level regularizer is
selected by
\begin{equation}
\lambda_c^{\star}=
\underset{\lambda\in\Lambda}{\operatorname*{arg\,max}}
\operatorname{median}_{j}\,R^{2}_{cj}(\lambda).
\label{eq:lambda_selection}
\end{equation}
Responses satisfying $R^{2}_{cj}(\lambda_c^{\star})\geq0.30$ form the monitored
set $\mathcal{M}_c$. If fewer than five responses pass this threshold, the best
$\max(5,\lfloor m_c/5\rfloor)$ responses are retained. This fallback guarantees
multichannel evidence while making weakly predictable channels less influential.
The out-of-fold residuals
$e^{\mathrm{cal}}_{ctj}=y_{ctj}-\widehat{y}^{(-k(t))}_{ctj}$ constitute the
calibration sample. The model is then refitted on all normal training rows and
applied once to $\mathcal{T}^{\mathrm{pr}}_c$ to obtain prediction residuals
$e^{\mathrm{pr}}_{ctj}$.

\subsection{Residual standardization and extrapolation-aware scaling}

The primary UC-RCF configuration estimates a channel scale from cross-fitted
residuals,
\begin{equation}
\widehat{\sigma}_{cj}=
\max\left\{\operatorname{sd}_{t}
\left(e^{\mathrm{cal}}_{ctj}\right),10^{-6}\right\}.
\label{eq:global_scale}
\end{equation}
To reduce overreaction when the operating vector lies outside well-supported
regions of the nonlinear design, this scale is adjusted by regularized
prediction leverage. For a new design vector $\boldsymbol{\phi}_{ct}$,
\begin{equation}
h_{ct}=\sum_{\ell}
\left(\boldsymbol{\phi}_{ct}^{\top}\mathbf{v}_{\ell}\right)^2
\frac{s_{\ell}^{2}}{(s_{\ell}^{2}+\lambda_c^{\star})^{2}},
\label{eq:regularized_leverage}
\end{equation}
and the nonconformity score becomes
\begin{equation}
z_{ctj}=
\frac{|e_{ctj}|}
{\widehat{\sigma}_{cj}\sqrt{1+h_{ct}}}.
\label{eq:nonconformity}
\end{equation}
Thus, an equally sized residual contributes less evidence when it occurs in a
poorly supported operating region. A heteroscedastic alternative that predicts
$\log(e_{ctj}^{2}+\epsilon)$ from the same design is retained as an ablation,
not silently combined with the primary configuration.

\subsection{Channel-wise conformal evidence}

For each monitored response $j$, let
$\mathcal{Z}_{cj}=\{z^{\mathrm{cal}}_{cij}\}_{i=1}^{n_{cj}}$ be its finite
cross-fitted nonconformity scores. The rank-based conformal value assigned to a
prediction score $z^{\mathrm{pr}}_{ctj}$ is
\begin{equation}
p_{ctj}=
\frac{1+\sum_{i=1}^{n_{cj}}
\mathbf{1}\!\left(z^{\mathrm{cal}}_{cij}\geq z^{\mathrm{pr}}_{ctj}\right)}
{n_{cj}+1}.
\label{eq:conformal_pvalue}
\end{equation}
Small values indicate that the response is unusual relative to its own
cross-fitted healthy behavior. For a calibration observation $i$, the analogous
leave-one-out rank is
\begin{equation*}
p^{\mathrm{cal}}_{cij}=
\frac{1+\sum_{r\neq i}
\mathbf{1}\!\left(z^{\mathrm{cal}}_{crj}\geq z^{\mathrm{cal}}_{cij}\right)}
{n_{cj}},
\end{equation*}
so that a score is never used to rank itself. Two aggregate statistics are
available:
\begin{align}
A^{\mathrm{mean}}_{ct}
&=\frac{1}{|\mathcal{M}_{ct}|}
\sum_{j\in\mathcal{M}_{ct}}-\log_{10}(p_{ctj}),
\label{eq:mean_evidence}\\
A^{\mathrm{frac}}_{ct}
&=\frac{1}{|\mathcal{M}_{ct}|}
\sum_{j\in\mathcal{M}_{ct}}
\mathbf{1}(p_{ctj}<\gamma),\qquad \gamma=0.01,
\label{eq:fraction_evidence}
\end{align}
where $\mathcal{M}_{ct}\subseteq\mathcal{M}_c$ contains responses observed at
time $t$. The primary cross-asset detector uses
$A_{ct}=A^{\mathrm{frac}}_{ct}$ because it measures the breadth of simultaneous
channel departures and is less dominated by a single extreme residual.

For every candidate quantile $q$, a case-specific point threshold is derived
from the corresponding out-of-fold aggregate scores,
\begin{equation}
\tau_c(q)=\operatorname{Quantile}_{q}
\left(\{A^{\mathrm{cal}}_{ct}\}_{t\in\mathcal{T}^{\mathrm{tr}}_c}\right),
\qquad
I_{ct}(q)=\mathbf{1}\!\left(A^{\mathrm{pr}}_{ct}>\tau_c(q)\right).
\label{eq:point_firing}
\end{equation}
This first calibration layer removes case-specific score units. The second,
cross-asset layer below selects $q$ according to the behavior of the complete
persistent-alarm process rather than pointwise exceedance alone.

\subsection{Persistence-sensitive criticality}

Isolated SCADA deviations should not trigger a case-level alarm. UC-RCF converts
the binary firing sequence into a reflected random walk. After retaining only
normal-status prediction rows in chronological order, define
\begin{equation}
d_{ct}(q)=2I_{ct}(q)-1,
\qquad
C_{ct}(q)=\max\{0,C_{c,t-1}(q)+d_{ct}(q)\},
\qquad C_{c0}=0.
\label{eq:criticality_recursion}
\end{equation}
The equivalent closed form is
$C_{ct}=S_{ct}-\min_{0\leq u\leq t}S_{cu}$, where
$S_{ct}=\sum_{v\leq t}d_{cv}$ and $S_{c0}=0$. A case alarm is declared when
\begin{equation}
\max_{t}C_{ct}(q)\geq H,\qquad H=72.
\label{eq:case_alarm}
\end{equation}
The $+1/-1$ update rewards sustained evidence and erases sporadic firings. The
counter is frozen when the operating-status filter is not satisfied, and no
post-alarm latching is applied.

\subsection{Cross-asset calibration of the alarm process}

Pointwise conformal ranks do not by themselves control the probability that a
long, dependent sequence crosses $H$. Moreover, the calibration and prediction
partitions may exhibit asset-dependent covariate shift. UC-RCF therefore
calibrates the complete criticality process using prediction-period sequences
from verified-normal donor cases.

Consider an outer fold that holds out asset $b$ and a target case $c$. The donor
set is
\begin{equation}
\mathcal{D}_{c,b}=
\left\{d\in\mathcal{C}: d\text{ is normal},\,
a(d)\neq b,\,a(d)\neq a(c)\right\}.
\label{eq:donor_set}
\end{equation}
The first exclusion protects the outer test asset; the second prevents a target
case from borrowing its own turbine's prediction behavior during calibration.
For each $q$ in
\begin{equation*}
\mathcal{Q}=\{0.70,0.80,0.85,0.90,0.925,0.95,0.965,0.98,0.99,0.995\},
\end{equation*}
the donor firing sequences are partitioned into windows with target-matched
length $L_c=\max(|\mathcal{T}^{\mathrm{pr,N}}_c|,200)$ and 144-sample stride.
If a donor is shorter than $L_c$, its full available sequence is used. Let
$M_{d\ell}(q)$ denote the maximum criticality in donor window $\ell$. The
empirical case-level null risk is
\begin{equation}
\widehat{\pi}_{c,b}(q)=
\frac{1}{N_{c,b}(q)}
\sum_{d\in\mathcal{D}_{c,b}}\sum_{\ell}
\mathbf{1}\!\left(M_{d\ell}(q)\geq H\right),
\label{eq:null_false_alarm}
\end{equation}
where $N_{c,b}(q)$ is the number of available donor windows. For a nominal
case-level false-alarm budget $\alpha$, UC-RCF selects the loosest admissible
point threshold,
\begin{equation}
q^{\star}_{c,b}(\alpha)=
\min\left\{q\in\mathcal{Q}:\widehat{\pi}_{c,b}(q)\leq\alpha\right\}.
\label{eq:q_selection}
\end{equation}
If no candidate satisfies the budget, $q=0.995$ is used and the case is marked
as calibration-infeasible. This flag is retained for the validity audit.

\subsection{Nested leave-one-asset-out policy selection}

The nominal budget is treated as an operating policy rather than tuned on a
held-out asset. For outer fold $b$, candidate values are
\begin{equation*}
\mathcal{A}=\{0.02,0.05,0.10,0.15,0.20,0.30,0.40,0.50,0.60\}.
\end{equation*}
For each $\alpha\in\mathcal{A}$, Eq.~\eqref{eq:q_selection} is recomputed with
asset $b$ excluded. The selected policy is
\begin{equation}
\alpha_b^{\star}=
\underset{\alpha\in\mathcal{A}}{\operatorname*{arg\,max}}
\;\operatorname{CARE}
\left(\mathcal{C}\setminus\{c:a(c)=b\};\alpha\right),
\label{eq:outer_policy}
\end{equation}
with deterministic grid-order tie breaking. Only then is
$\alpha_b^{\star}$ applied to cases belonging to $b$. The final endpoint pools
the predictions produced when each asset is held out. Any detector variant or
residual-scale configuration selected from labels must be nested in the same
outer loop; alternatively, it must be fixed before outer evaluation. This rule
is essential because asset-disjoint policy selection alone cannot correct a
globally outcome-selected model configuration.

The full procedure is summarized as follows:
\begin{enumerate}
    \item For each case, construct Eq.~\eqref{eq:nonlinear_design} using normal
    training statistics only.
    \item Generate embargoed out-of-fold residuals, select
    $\lambda_c^{\star}$ and $\mathcal{M}_c$, and refit the multi-output model on
    all normal training rows.
    \item Compute Eqs.~\eqref{eq:regularized_leverage}--\eqref{eq:fraction_evidence}
    on the calibration and prediction partitions.
    \item For each outer held-out asset and candidate $\alpha$, build donor
    null windows, select $q^{\star}_{c,b}(\alpha)$, and form the criticality
    trajectory for every case.
    \item Select $\alpha_b^{\star}$ using development assets only, apply it to
    the held-out asset, and pool all held-out predictions exactly once.
\end{enumerate}

\begin{table}[htbp]
\centering
\small
\caption{Fixed UC-RCF settings used in the proposed method.}
\label{tab:ucrcf_settings}
\begin{adjustbox}{max width=\textwidth}
\begin{tabular}{lll}
\toprule
Component & Setting & Interpretation \\
\midrule
Input statistic & average channels only & common SCADA statistic \\
Normal status IDs & $\{0,2\}$ & model-fitting and alarm-update filter \\
Cross-fitting & 5 folds; 432-sample blocks & 3-day contiguous blocks \\
Embargo & 144 samples & 24~h around held blocks \\
Rolling windows & $\{6,36,144\}$ & 1, 6, and 24~h context \\
Core drivers & at most 6 & interactions and rolling terms \\
Minimum fit size & 3000 rows & case inclusion requirement \\
Target screen & out-of-fold $R^2\geq0.30$ & minimum predictability \\
Channel level $\gamma$ & 0.01 & fraction-evidence statistic \\
Criticality threshold $H$ & 72 & persistent-alarm boundary \\
Null-window stride & 144 samples & 24~h between window starts \\
Post-alarm latching & disabled & trajectory remains observable \\
\bottomrule
\end{tabular}
\end{adjustbox}
\end{table}

\subsection{Interpretation and scope of calibration}

The rank construction in Eq.~\eqref{eq:conformal_pvalue} is distribution-free
only under exchangeability of the standardized residuals used for calibration
and prediction. Serial dependence and cross-asset distribution shift prevent an
automatic finite-sample guarantee at the case-alarm level. UC-RCF does not hide
this limitation behind the term ``conformal.'' Instead, it uses asset-disjoint
prediction-period null sequences to target the operational false-alarm budget
and reports both the observed false-alarm rate and its ratio to the nominal
budget. Consequently, the method separates three questions that are often
conflated: whether a channel residual is unusual, whether unusual evidence is
persistent, and whether the resulting alarm process transfers to a different
turbine.
